@@M0@@
\item \uline{The cokernel} @@M1@@ \uline{of} @@M2@@ \uline{is a locally free} @@M3@@\uline{-module.}
\item \uline{The homomorphism} @@M4@@ \uline{is surjective.}
\end{enumerate}
\uline{Then the canonical morphism} @@M5@@ \uline{is representable by an affine morphism of finite presentation; more precisely, there is a canonical exact sequence of} f.p.p.f. \uline{sheaves,}
@@M6@@
\uline{and, for} @@M7@@\uline{, the monomorphism} @@M8@@ \uline{is representable by a closed immersion and} @@M9@@ \uline{is a torsor over} @@M10@@ \uline{under the group} @@M11@@ \uline{(that is, the sequence of} f.p.p.f. \uline{sheaves}
@@M12@@
\uline{is exact).}

For an @@M13@@-scheme @@M14@@, consider the following two diagrams with exact rows:
@@M15@@
@@M16@@

Observe carefully that the square on the left of this diagram \uline{is not necessarily commutative}. Nevertheless, we have:
\sgaStableTarget{sga6.e12.para.3-2}@@M17@@

